

\newpage

%\footline={\hss\folio\hss}
%\pageno=1

\dspace
\subhead
6. Profondeur 
\endsubhead

\subhead
6.1
\endsubhead
Avec les notations de 5.1, pour tout point
base $x$, l'inclusion de $X$ dans
$A^*$ induit un morphisme du
groupe fondamental motivique de $X$ dans
celui, $\dbQ(1)$, de $A^*$.
On d\'efinit la filtration par la profondeur
$D$ de $\pi_1(X,x)$ par
$$
\align
&D^0\pi_1(X,x):=\pi_1(X,x)\tag6.1.1\\
&D^1\pi_1(X,x):=\Ker(\pi_1(X,x)
\to \pi_1(A^*,x)
\endalign
$$
et, pour $n\Ge 1$, 
$$
\align
D^n\pi_1(X,x):=&\text{le $n^{\ieme}$ groupe de la}\\
&\text{s\'erie centrale descendante de
$D^1\pi_1(X,x)$}.
\endalign
$$
On filtre de m\^{e}me les alg\`ebres de Lie:
$$
D^n\Lie\,\pi_1(X,x)=\Lie\,D^n\pi_1(X,x).
\tag6.1.2
$$

Si on pousse le $\pi_1(X,x)$-torseur $P_{y,x}$
par $\pi_1\to \pi_1/D^n$, les
$P_{y,x}/D^n$ obtenus forment un
groupo{\umi}de quotient du groupo{\umi}de
des $P_{y,x}$.

Appliquons A.10, A.11.
La filtration par la profondeur induit une
filtration de l'alg\`ebre enveloppante
compl\'et\'ee $\scrU\uphat\Lie\,\pi_1(X,x)$
ainsi que de l'alg\`ebre affine
$\scrO(\pi_1(X,x))$, qui en est duale.
Elle induit aussi  une filtration sur les
alg\`ebres affines des torseurs $P_{y,x}$.
L'alg\`ebre affine du quotient
$\pi_1(X,x)/D^{n+1}$ de $\pi_1(X,x)$ est la
sous-alg\`ebre de l'alg\`ebre affine de
$\pi_1(X,x)$ engendr\'ee par $D^{-n}
\scrO(\pi_1(X,x))$.
De m\^{e}me pour les $P_{y,x}$.

\subhead
6.2
\endsubhead
En r\'ealisation $\omega$, l'alg\`ebre de Lie
de $\pi_1(X,x)$ est $\scrL^{\uphat}$ (5.8), le
compl\'et\'e de l'alg\`ebre de Lie libre
$\scrL$ engendr\'ee par $e_0$ et les $e_\zeta$
($\zeta\in\mu_N$).
La filtration de profondeur 
est d\'eduite 
la graduation par le degr\'e en les $e_\zeta$,
appel\'e le {\it $D$-degr\'e}.
Cette graduation 
est compatible \`a la
graduation provenant de celle de $\omega$ pour
laquelle $e_0$ et les $e_\zeta$ sont de
degr\'e un et est $\mu_N$-\'equivariante.
Elle se propage aux alg\`ebres
enveloppantes, aux alg\`ebres affines et aux
alg\`ebres affines des torseurs triviaux
$(P_{y,x})_\omega$.
L'action correspondante $\rho_D$ de $\dbG_m$
pour laquelle $\rho_D(\lambda)$ fixe $e_0$ et
transforme $e_\zeta$ en $\lambda e_\zeta$ sera
appell\'ee action ``$D$-degr\'e''.

On prendra garde que les filtrations $D$ sont
motiviques et stables par l'automorphisme
$u\mapsto u^{-1}$ de $X$, mais que le
$D$-degr\'e est d\'efini seulement en
r\'ealisation $\omega$ et n'est pas stable par
$u\mapsto u^{-1}$.
L'automorphisme de $\scrL\uphat$ induit par
$u\mapsto u^{-1}$  \'echange $e_0$, qui est
homog\`ene de $D$-degr\'e $0$, et
$e_\infty=-e_0-\sum e_\zeta$,  qui ne l'est
pas.

\subhead
6.3
\endsubhead
Le $\MT(\scrO[1/N])$-sch\'ema en groupe $V$
de 5.15 agit sur $P_{0,0}$ et $P_{1,0}$.
On d\'efinit la filtration $D$ de $\Lie\,V$
par la condition que $D^p\Lie\,V$ soit maximal
tel que l'action de $\Lie\,V$ sur
$\scrO(P_{0,0})$ et $\scrO(P_{0,1})$ envoie
$D^p\Lie\,V\otimes D^q\scrO(P_{0,0})$ dans
$D^{p+q}\scrO(P_{0,0})$ et $D^p\Lie\,V\otimes
D^q\scrO(P_{1,0})$ dans
$D^{p+q}\scrO(P_{1,0})$.

Passons \`a la r\'ealisation $\omega$, o\`u
$V_\omega$ s'identifie \`a
$\left(\prod,\circ\right)$ agissant sur les copies
$\prod_0$ et $\prod_{1,0}$ de $\prod$ par
$g\mapsto\left<g\right>_0$ et
$\left<g\right>_{1,0}$.

\proclaim{Proposition 6.4}
L'action ``$D$-degr\'e'' $\rho_D$ de $\dbG_m$
sur $\prod$ respecte les actions
$\left<g\right>_0$ et $\left<g\right>_{10}$ de
$\left(\prod,\circ\right)$ sur $\prod$.
Elle respecte donc la loi de groupe $\circ$ sur
$\prod$ et le crochet $\{\,\,,\,\,\}$ de
$\scrL\uphat$.
\endproclaim

L'action $\left<g\right>_0$ de $g\in\prod$ sur
le groupe $\prod$ est caract\'eris\'ee par les
propri\'et\'es de commuter \`a l'action de
$\mu_N$ sur $\prod$, de fixer $\exp(e_0)$ et
d'envoyer $\exp(e_1)$ sur $g^{-1}\exp(e_1)g$.
Son conjugu\'e par $\rho_D(\lambda)$ commute
encore \`a $\mu_N$, fixe $\exp(e_0)$ et envoie
$\exp(\lambda e_1)$ sur\break
$\rho_D(\lambda)(g)^{-1}\exp(\lambda e_1)
\rho_D(\lambda)(g)$.
Cette derni\`ere propri\'et\'e \'equivaut \`a
envoyer $\exp(e_1)$ sur
$\rho_D(\lambda)(g)^{-1}\exp(e_1)
\rho_D(\lambda)(g)$: $\rho_D(\lambda)$
conjugue $\left<g\right>_0$ en
$\left<\rho_D(\lambda)(g)\right>_0$.
Conjugant $\left<g\right>_{10}\colon\,x\mapsto
g\left<g\right>_0(x)$ par $\rho_D(\lambda)$,
on en d\'eduit que $\rho_D(\lambda)$
transforme $\left<g\right>_{10}$ en
$\left<\rho_D(\lambda)(g)\right>_{10}$.

\subhead
6.5
\endsubhead
Puisque l'action de $\Lie\,V_\omega=(\scrL\uphat,
\{\,\,,\,\,\})$ sur $\scrO(P_{1,0})_\omega$
est fid\`ele, il r\'esulte de 6.4 que la
filtration $D$ de $\Lie\,V_\omega$ co{\umi}ncide
avec la filtration $D$ de $\scrL\uphat$
d\'eduite de la graduation par le $D$-degr\'e
de $\scrL$.
En particulier, $D^0\Lie\,V_\omega=\Lie\,V_\omega$.

Transportons \`a $V_\omega$ l'action
``$D$-degr\'e'' $\rho_D$ de $\dbG_m$ sur
$\left(\prod,\circ\right)$.
Cette action gradue $\Lie(V_\omega)$, son
alg\`ebre enveloppante universelle
compl\'et\'ee et l'alg\`ebre affine
$\scrO(V_\omega)$.
Les filtrations correspondantes
co{\umi}ncident avec celles d\'efinie par la
filtration $D$ de $\Lie\,V_\omega$ par les
$D^p\Lie\,V_\omega$.
Il en r\'esulte que l'isomorphisme de 
sch\'emas de $V_\omega$ 
avec $\prod$ induit un isomorphisme de
$V_\omega/D^pV_\omega$ avec $\prod/D^p\prod$,
les alg\`ebres affines des quotients \'etant
dans les deux cas les sous-alg\`ebres
engendr\'ees par $D^{-p}$ des alg\`ebres
affines (A.10).
Plus pr\'ecis\'ement, l'action de $V_\omega$
sur $\prod_{10}$ fait de $\prod_{10}$ un
espace principal homog\`ene sous $V_\omega$,
trivialis\'e par $1_{10}$ fixe sous $\rho_D$,
et induit une action principale homog\`ene de
$V_\omega/D^pV_\omega$ sur $\prod_{10}/D^p$.
Ind\'ependamment du foncteur fibre, il reste
vrai que l'action de $V$ sur $P_{0,0}$ et
$P_{1,0}$ induit une action de $V/D^pV$ sur
$P_{0,0}/D^p$ et $P_{1,0}/D^p$, et que
$P_{1,0}/D^p$ est un espace principal
homog\`ene sous $V/D^pV$. 

\proclaim{Proposition 6.6}
Le morphisme 5.15 de $U$ dans $V$ se factorise
par $D^1V$.
\endproclaim

Le quotient $D^1V/D^2$ \'etant ab\'elien, ce
morphisme induit
$$
U^{\ab}\to D^1V/D^2V.\tag6.6.1
$$

\demo{Preuve}
D'apr\`es 6.5, il suffit de v\'erifier que $U$
agit trivialement sur $P_{1,0}/D^1$, le
$\dbQ(1)$-torseur des chemins, dans
$A^*$, de $1_0$ \`a $1$.
En effet, $P_{\mu,\lambda_0}$
($A^*$) est canoniquement
isomorphe \`a
$P_{\mu,\lambda}(A^*)$, et,
faisant $\lambda=\mu$, on voit que
$P_{1,0}/D^1$ est le $\dbQ(1)$-torseur trivial,
sur lequel $U$ agit trivialement.
\enddemo

\subhead
6.7
\endsubhead
Nous nous proposons de montrer que (6.6.1) est
injectif et de calculer son image.
Noter que l'action de $U$ sur $U^{\ab}$, ainsi
que l'action de $U$ sur $D^1V/D^2V$
(d\'eduite de l'action int\'erieure de $V$ sur
lui-m\^{e}me), sont triviales.
Les deux membres de (6.6.1) sont donc des
produits de sch\'emas en groupe vectoriels
$\dbQ(n)$, et le passage \`a la r\'ealisation
$\omega$ ne perd aucune information.

En r\'ealisation $\omega$, le noyau $D^1\scrL$
de $\Lib(e_0,(e_\zeta))\to\Lib(e_0)$:
$e_0\mapsto e_0$, $e_\zeta\mapsto 0$, est,
pour le crochet [\quad], l'alg\`ebre de $\Lie$
libre engendr\'ee par les $(\ad\,e_0)^n(e_\zeta)$
($n\Ge 0$, $\zeta\in\mu_N)$, de $D$-degr\'e
$1$.
Le quotient $D^1\scrL/D^2\scrL$ a donc pour
base les images des $(\ad\,e_0)^n(e_\zeta)$,
$\ad\,e_0$ \'etant pris au sens du crochet
$[\,\,,\,\,]$.

Notons $\scrE^n$ la composante de poids $-2n$
de $D^1\scrL/D^2\scrL$, de base les images des\break
$(\ad\,e_0)^{n-1}(e_\zeta)$.
Le quotient $\Gr_D^1(\scrL\uphat)$ est le
produit des $\scrE^n$.
Quand $n$ sera fix\'e, nous \'ecrisons
simplement $\scrE$ pour $\scrE^n$, et
$E_\zeta$ pour la base des
$(\ad\,e_0)^{n-1}(e_\zeta)$.

Fixons un plongement complexe $\sigma$ de $k$.
L'\'el\'ement $\dch(\sigma)$ de
$V_\omega(\dbC)\Subset\dbC\ll e_0,(e_\zeta)\gg$
a un coefficient de $e_0$ nul (5.17).
Il est donc dans $D^1V_\omega(\dbC)$.
Calculons sa projection dans
$D^1V_\omega/D^2V_\omega$, d'alg\`ebre de Lie
$D^1\scrL\uphat/D^2\scrL\uphat$.
Pour
$$
x=\sum\lambda_{n,\zeta}(\ad e_0)^{n-1}(e_\zeta)
+\text{\rm termes de $D$-degr\'e $\Ge 2$},
$$
l'exponentielle de $x$ dans $V_\omega=\prod$
est donn\'ee par 5.14 (ii).
On a
$$
\exp(x)=1+x+\text{\rm termes de $D$-degr\'e $\Ge
2$},
$$
et $\lambda_{n,\zeta}$ est donc le coefficient
de $e_0^{n-1}e_\zeta$ dans $x$.

D'apr\`es 5.17, la projection de
$\dch(\sigma)$ dans $D^1V_\omega/D^2V_\omega$
est donc l'exponentielle de l'\'el\'ement
$$
\dch^1(\sigma)=\sum\limits_{n,\zeta}
\biggl(\sum\limits_{m}
\frac{\sigma(\zeta^{-m})}{m^n}\biggr)
(\ad\,e_0)^{n-1}(e_\zeta)
$$
de son alg\`ebre de Lie $\prod\scrE^n$.
Pour $n=1$, $\zeta=1$, la somme sur $m$ diverge
et est \`a remplacer par $0$.

Identifions $U^{\ab}$ et $D^1V/D^2V$ \`a
leurs alg\`ebre de Lie par l'application
exponentielle et projetons sur $U^{\ab}$ et
$D^1V/D^2V$.
On obtient qu'il existe $a_\sigma$ dans 
$\gru_\omega^{\ab}
\widehat{\otimes}\dbC$ v\'erifiant
$$
\align
&\text{Image de $a_\sigma$ dans
$\left(\Lie\,D^1V_\omega/D^2V_\omega=\prod \scrE^n\right)\widehat{\otimes}\dbC $}
\tag6.7.1\\
&=\text{$\dch^1(\sigma)+$ \'el\'ement de
$\prod(2\pi i)^n\scrE^n$}.
\endalign
$$

Notons  $\scrI$ l'image de Lie $U_\omega$ dans Lie
$V_\omega$.
La filtration $D$ de Lie $V$ induit une
filtration $D$ de $\scrI$.
On prendra garde que l'action ``$D$-degr\'e''
de $\dbG_m$ sur $\Lie\,V_\omega$ n'a aucune
raison de respecter $\scrI$ (et en fait ne
respecte pas $\scrI$).

\proclaim{Th\'eor\`eme 6.8}\enspace
{\rm (i)} Le morphisme (6.6.1) est injectif; il
induit donc un isomorphisme de
$\gru_\omega^{\ab}$ avec $\scrI/D^2\scrI$.

\noindent
{\rm (ii)} En r\'ealisation $\omega$ et en poids
$-2n$, l'image dans $\scrE:=\scrE^n$ de
$\omega_n(\gru^{\ab})$ est le sous-espace
$\scrE_{\distr}^+$ des $\sum x_\zeta E_\zeta$
v\'erifiant la relation de parit\'e
$$
x_{\zeta^{-1}}=(-1)^{n-1}x_\zeta\tag6.8.1
$$
et la relation de distribution de poids $n-1$
(resp. de distribution trou\'ee si $n=1$): si
$N=Md$ avec $d\Ge 1$ et que $\eta^M=1$, on a
$$
x_\eta=\frac{1}{d^{n-1}}\sum\limits_{\zeta^d=\eta}
x_\zeta\tag6.8.2
$$
(resp. $x_1=0$ et (6.8.2) seulement pour
$\eta\not=1$ si $n=1$).
\endproclaim

\remark\nofrills{\bf Remarque 6.9 {\rm (i)}.}\enspace
La relation de parit\'e (6.8.1) est la
relation (6.8.2) pour $d=-1$.

(ii)\enspace
Le $\dbQ$-espace vectoriel $\scrE_{\distr}^+$
est muni de formes lin\'eaires $x\mapsto
x_\zeta$ v\'erifiant les relations de parit\'e
et de distribution de 6.8 (ii), et est
universel pour cette propri\'et\'e.
Son dual $\scrD_{\distr}^+$ est muni d'une
application $\mu_N\to\scrD_{\distr}^+$ v\'erifiant
les m\^{e}mes identit\'es, et est universel
pour cette propri\'et\'e: c'est l'espace dans
lequel prend ses valeurs la distribution de
poids $n-1$ (resp. trou\'ee si $n=1$) de
parit\'e $(-1)^{n-1}$ universelle.
\endremark

\demo{6.10 Preuve de 6.8}
Pour $\sigma$ un plongement complexe de $k$,
posons
$$
d_\sigma:=\sum\limits_{\zeta}\sum\limits_{m}
\frac{\sigma(\zeta^{-m})}{m^n}E_\zeta.
\tag6.10.1
$$
Pour $n=1$, le coefficient, divergent, de
$E_1$ est \`a remplacer par z\'ero.

D'apr\`es 6.7.1, il existe $a_\sigma$ dans
l'image de $\omega_n(\gru^{\ab})$ tel que
$$
d_\sigma\equiv a_\sigma\pmod{(2\pi i)^n\scrE}.
\tag6.10.2
$$
\enddemo

\medskip\noindent
{\sans Cas $n=1$.}\enspace
Fixons $\sigma$ et identifions $k$ \`a son
image dans $\dbC$ par $\sigma$.
En particulier, \'ecrivons $\zeta$ pour
$\sigma\zeta$.
Avec cette notation, (6.10.1) se r\'ecrit
$$
(d_\sigma)_\zeta=-\log(1-\zeta).
$$

D'apr\`es (2.3.1), $\omega_n(\gru^{\ab})$ est
le dual de $\scrO_N[1/N]^*\otimes\dbQ$.
Si $\scrO_N^+$ est la partie r\'eelle de
$\scrO_N$, c'est encore le dual de
$\scrO_N^+[1/N]^*\otimes\dbQ$.
Soit $\scrF$ son image dans $\scrE$.

Puisque $a_\sigma$ est dans $\scrF_\dbC$, sa
partie r\'eelle est dans $\scrF_\dbR$.
Prenant la partie r\'eelle de (6.10.2), on
obtient que
$$
-\sum \log\vert1-\zeta\vert~E_\zeta
$$
est dans $\scrF_\dbR$.

D'apr\`es un th\'eor\`eme de Bass (1966) (cf.
Washington (1997) 8.9 et 12.18),
$\scrO_N[1/N]^*\otimes\dbQ$, muni de
$\zeta\mapsto 1-\zeta$ (resp. $1$ si
$\zeta=1$) est la distribution trou\'ee paire
universelle sur $\mu_N$.
Puisque $\scrO_M^+[1/N]^*
\otimes\dbQ\simover\scrO_N[1/N]^*\otimes\dbQ$, 
l'application 
$x\mapsto \log\,\vert x\vert$ est injective sur
$\scrO_N[1/N]^*/\torsion$ et il n'y a donc
d'autres relations rationnelles entre les
$\log\,\vert1-\zeta\vert$ que celles
d\'eduites de (6.8.1) et (6.8.2).
L'image  $\scrF$ contient donc
$\scrE_{\distr}^+$.
Puisque $\scrE_{\distr}^+$ a le m\^{e}me rang
que $\scrO_N^*[1/N]^*\otimes\dbQ$, (i) et (ii)
en r\'esultent.

\medskip\noindent
{\sans Cas $n>1$.}\enspace
Soit $\scrF$ l'image de
$\omega_n*(\gru^{\ab})$ dans $\scrE$.
Prenant, selon la parit\'e de $n$, la partie
r\'eelle ou la partie imaginaire de (6.10.2),
on obtient que pour tout $\sigma$
$$
d_\sigma+(-1)^{n-1}\dbar_\sigma
$$
est dans $\scrF_\dbC$.
Il r\'esulte de (6.10.1) et des identit\'es
$$
\align
&\sum\limits_{\zeta^d=\eta}\zeta^m
=\cases
0 &\text{si $d\nmid m$}\\
d\eta^{m/d} &\text{si $d\mid m$}
\endcases\\
\vspace{5pt}
&\zeta^{-1}=\zetabar
\endalign
$$
que les $d_\sigma+(-1)^{n-1}\dbar_\sigma$ sont
dans $(\scrE_{\distr}^+)_\dbC$.

\proclaim{Lemme 6.11}
Le sous-espace $W$ de $\scrE_\dbC$ engendr\'e
par les $d_\sigma\pm (-1)^{n-1}\dbar_\sigma$
est de dimension au moins \'egale au nombre de
caract\`eres de $(\dbZ/N\dbZ)^*$ de parit\'e
$(-1)^{n-1}$.
\endproclaim

D\'eduisons 6.8 de 6.11.
La borne inf\'erieure 6.11 pour la dimension
de $W$ co{\umi}ncide tant avec la dimension de
$\omega(\gru^{\ab})_n$ (2.3.1) qu'avec la
dimension de $\scrE_{\distr}^+$ (Kubert (1979)
App.).
Puisque $W$ est contenu dans l'image
complexifi\'ee $\scrF_\dbC$, $W$ est \'egal
\`a cette image, (6.1.1) est injectif.
De m\^{e}me, puisque $S$ est contenu dans
$(\scrE_{\distr}^+)_{\dbC}$, on a \'egalit\'e,
prouvant 6.8.

\demo{Preuve de $6.11$}
On peut supposer que $k=\dbQ(\exp(2\pi
i/N))\Subset \dbC$.
On identifie l'ensemble des plongements
complexes de $k$ \`a $(\dbZ/N)^*$ par
$\zeta\mapsto \zeta^a$, et $\mu_N$ \`a
$\dbZ/N$ par $b\mapsto\exp(2\pi i\,b/N)$.
Avec ces notations, $W$ est l'image de
l'application
$$
F\colon\, \dbC^{(\dbZ/N)^*}\to
\dbC^{\dbZ/N}
$$
de matrice
$$
\sum\limits_{m}\,
\frac{\exp(-2\pi i\,abm)}{m^n}
+(-1)^{n-1}
\frac{\exp(-2\pi i\,abm)}{m^n}\,\,.
$$
Il nous faut minorer le rang de $F$.
Il sera plus commode de consid\'erer
l'application transpos\'ee ${}^tF$.
Les applications $F$ et ${}^tF$ sont
\'equivariantes, pour les actions $e_x\mapsto
e_{ax}^{-1}$ et $e_x\mapsto e_{ax}$ de
$(\dbZ/N)^*$ sur les deux membres, et il
suffit de v\'erifier que pour tout caract\`ere
$\chi$ de $(\dbZ/N)^*$ de parit\'e
$(-1)^{n-1}$, la $\chi$-partie de ${}^tF$
est non nulle.
Soient $M$ le diviseur de $N$ tel que $\chi$
provienne d'un caract\`ere primitif $\chi_0$
de $(\dbZ/M)^*$, et prolongeant $\chi_0$ par $0$
sur $\dbZ/M$.
V\'erifions que l'image par ${}^tF$ du vecteur
$(\chi_0(b))_{b\in\dbZ/N}$ est non nulle.
Prenant la composante $a=1$, on se ram\`ene
\`a v\'erifier que
$$
\sum\limits_{b,m}\,
\frac{\exp(2\pi i\,bm).\chi_0(b)}{m^n}
\not=0
$$
C'est en effet $N/M$ fois le produit de la
somme de Gau{\ss} $\sum\exp(2\pi
i\,b)\chi_0(b)$ (somme sur $(\dbZ/M)^*$)
par $L(\chi_0,n)$.
\enddemo

\remark{\bf Remarque 6.12}
L'argument utilis\'e pour $n>1$, utilisant
tous les plongements complexes et (6.10.2)
modulo $(2\pi i)^n\scrE\otimes\dbR$
plut\^{o}t qu'un seul plongement et (6.10.2)
modulo $2\pi i\scrE$, s'applique aussi pour
$n=1$ si $N$ n'a qu'un seul facteur premier.
Dans ce cas en effet, l'application
r\'egulateur de $\scrO_N^*[1/N]$ dans $\dbR^m$,
$m$ le nombre de places \`a l'infini, a un
noyau fini et  une image discr\`ete.
\endremark

\remark{\bf Remarque 6.13}
Combinant les r\'esultats de cet article avec
ceux de Goncharov (2001), on obtient
imm\'ediatement l'analogue motivique de
r\'esultats de loc. cit. sur les alg\`ebres de
Lie de Galois: remplacer syst\'ematiquement
l'alg\`ebre de Lie de Galois $g_N^\ell$
par $\im(\Lie\,U_\omega\to\Lie\,V_\omega)$
(voir 10.5.3) qui est son analogue motivique,
et travailler sur $\dbQ$ plut\^{o}t que
$\dbQ_\ell$.
\endremark

